% IEEE numbered reference style.
% With thebibliography, the printed number of an entry is its POSITION in this list.
% Entries are in order of first citation, and uncited entries were removed (10 Oct 2026; again 11 Oct 2026 after the length cut).
% A new source goes where it is first cited; the text uses \cite, so numbers update.

\begin{thebibliography}{99}
% Verified 2026-09-24 via Crossref and Semantic Scholar abstract: uses mean
% square of unsprung-mass acceleration divided by vehicle speed.
% Camera and ML papers (11 Oct 2026): details from Report/references_pdf/camera_ml/README.md, which checked
% each against arXiv and Crossref; PDFs in that folder.
% Published ECCV author list (8 authors); the arXiv version lists 10.
% Published ICCV author list (2 authors); the arXiv version lists 6.
% Road-obstacle and anomaly detection, cited in the Literature Review (moved from Recommendations, 11 Oct 2026).
% NEW - added for the VL53L5CX proximity-alerting subsection (Methodology,
% Section on non-ML alternatives). Open-access; confirmed via PMC/MDPI.
% Verified 2026-09-24: genuine 652-page manual, documents every AT command
% used in Ch.3. Its +CGNSSINFO format (16 fields) does NOT match the 9-field
% response the fitted modem returned, so this is the version consulted, not
% confirmed as the version matching the module's firmware (see Ch.3 GNSS).
% Verified 2026-09-28: manual build dated 2026-09-15, p. 21 "Image Formats":
% "RGB888 - 24 bits per pixel, ordered [B, G, R]."
% Software and datasets (11 Oct 2026). VERIFY: access dates, and for the Roboflow set the exact title on its page.
% VERIFY: author order and that this SATC 2015 paper is the one the dataset's owners ask to be cited
% (the training README names the collectors as Nienaber, Kroon and Booysen, Stellenbosch University).
\bibitem{unwup2019} United Nations, Department of Economic and Social Affairs, Population Division, \emph{World Urbanization Prospects: The 2018 Revision}. New York, NY, USA: United Nations, 2019.

\bibitem{jinpengmanual} Jiangsu Jinpeng Group Co., Ltd, \emph{Electric Tricycle User Manual}. Vehicle documentation, TUKZIE Rev 0 platform.

\bibitem{adamsrasara2026} N. Adams and S. Rasara, ``Jinpeng Electric Tricycle Drivetrain: Characterization \& Modelling,'' TUKZIE vacation work report, University of Cape Town, Department of Electrical Engineering, Jul. 2026.

\bibitem{who2023} World Health Organization, \emph{Global Status Report on Road Safety 2023: Summary}. Geneva, Switzerland: World Health Organization, 2023.

\bibitem{nhtsa2013} National Highway Traffic Safety Administration, \emph{Visual-Manual NHTSA Driver Distraction Guidelines for In-Vehicle Electronic Devices}, Federal Register, vol. 78, no. 81, pp. 24818--24890, Apr. 2013.

\bibitem{winberg2026} S. Winberg and S. Jayalath, ``Development of an Embedded HPEC Telemetry Unit and Windshield HUD for the TUKZIE Rev 0 Platform,'' EEE4022S project brief, University of Cape Town, Department of Electrical Engineering, 2026.

\bibitem{zhu2021} Y. Zhu, Y. Jing, M. Jiang, Z. Zhang, D. Wang, and W. Liu, ``An experimental study of the cognitive load of in-vehicle multiscreen connected HUD,'' in \emph{Design, User Experience, and Usability: Design for Contemporary Technological Environments}, Lecture Notes in Computer Science, vol. 12781. Cham, Switzerland: Springer, 2021, pp. 268--285, doi: 10.1007/978-3-030-78227-6\_20.

\bibitem{kimmel2020} W. M. Kimmel, C. D. Weiss, and A. M. Mazzuchi, \emph{Technology Readiness Assessment Best Practices Guide}, NASA/SP-20205003605, 2020.

\bibitem{sebwa2025} O. S. Sebwa, P. V. Chombo, G. B. Gerutu, R. O. Kivugo, J. H. Hosea, and K. A. Greyson, ``Energy, economy and environmental assessment of e-tricycle for medium-size cargo deliveries: A case of Dar es Salaam City,'' \emph{The Journal of Engineering}, vol. 2025, no. 1, Art. no. e70073, 2025, doi: 10.1049/tje2.70073.

\bibitem{saenz2016} J. Saenz, M. Figliozzi, and J. Faulin, ``Assessment of the carbon footprint reductions of tricycle logistics services,'' \emph{Transportation Research Record}, vol. 2570, no. 1, pp. 48--56, 2016, doi: 10.3141/2570-06.

\bibitem{fhwa2016} Federal Highway Administration, \emph{Measuring and Specifying Pavement Smoothness}, Tech. Brief FHWA-HIF-16-032, Washington, DC, USA, 2016.

\bibitem{gonzalez2008} A. Gonz\'alez, E. J. O'Brien, Y.-Y. Li, and K. Cashell, ``The use of vehicle acceleration measurements to estimate road roughness,'' \emph{Vehicle System Dynamics}, vol. 46, no. 6, pp. 483--499, 2008, doi: 10.1080/00423110701485050.

\bibitem{sattar2018} S. Sattar, S. Li, and M. Chapman, ``Road surface monitoring using smartphone sensors: A review,'' \emph{Sensors}, vol. 18, no. 11, Art. no. 3845, 2018, doi: 10.3390/s18113845.

\bibitem{alatoom2025} Y. I. Alatoom and O. Smadi, ``Embedded framework for low-cost pavement condition evaluation using microcontroller and single-board computer platforms,'' \emph{Automation in Construction}, vol. 178, Art. no. 106442, 2025, doi: 10.1016/j.autcon.2025.106442.

\bibitem{bui2026} V.-H. Bui, D.-N. Tran, P. Q. Huy, N. V. Cong, and T.-H. Dao, ``Signal processing-driven real-time rough road detection with GPS accuracy improvement on edge microcontrollers,'' in \emph{Proc. IEEE 15th Int. Conf. Communication Systems and Network Technologies}, 2026, doi: 10.1109/CSNT69054.2026.11502079.

\bibitem{yiliguoqi2026} B. Yiliguoqi, B. Siriguleng, Arong, G. Guo, and J. Wang, ``Research on the evaluation and analysis of road surface roughness based on smartphone sensors and SVM,'' \emph{Scientific Reports}, vol. 16, Art. no. 4409, 2026, doi: 10.1038/s41598-025-34396-3.

\bibitem{li2018} Z. Li, W. Yu, and X. Cui, ``Online Classification of Road Roughness Conditions with Vehicle Unsprung Mass Acceleration by Sliding Time Window,'' \emph{Shock and Vibration}, vol. 2018, Art. no. 5131434, 2018, doi: 10.1155/2018/5131434.

\bibitem{alessandroni2017} G. Alessandroni, A. Carini, E. Lattanzi, V. Freschi, and A. Bogliolo, ``A study on the influence of speed on road roughness sensing: The SmartRoadSense case,'' \emph{Sensors}, vol. 17, no. 2, Art. no. 305, 2017, doi: 10.3390/s17020305.

\bibitem{iso2631} International Organization for Standardization, \emph{Mechanical Vibration and Shock: Evaluation of Human Exposure to Whole-Body Vibration, Part 1: General Requirements}, ISO 2631-1:1997, including Amendment 1:2010.

\bibitem{ru2022} X. Ru, N. Gu, H. Shang, and H. Zhang, ``MEMS inertial sensor calibration technology: Current status and future trends,'' \emph{Micromachines}, vol. 13, no. 6, Art. no. 879, 2022, doi: 10.3390/mi13060879.

\bibitem{tedaldi2014} D. Tedaldi, A. Pretto, and E. Menegatti, ``A robust and easy to implement method for IMU calibration without external equipments,'' in \emph{Proc. IEEE Int. Conf. Robotics and Automation (ICRA)}, 2014, pp. 3042--3049, doi: 10.1109/ICRA.2014.6907297.

\bibitem{webering2022} F. Webering, S. Kleinjohann, N. Stanislawski, and H. Blume, ``Improved calibration procedure for wireless inertial measurement units without precision equipment,'' arXiv:2207.04801 [eess.SP], 2022.

\bibitem{depasquale2010} G. De Pasquale and A. Som\`a, ``Reliability testing procedure for MEMS IMUs applied to vibrating environments,'' \emph{Sensors}, vol. 10, no. 1, pp. 456--474, 2010, doi: 10.3390/s100100456.

\bibitem{rfc5905} D. Mills, J. Martin, Ed., J. Burbank, and W. Kasch, ``Network Time Protocol Version 4: Protocol and algorithms specification,'' Internet Engineering Task Force, RFC 5905, Jun. 2010, doi: 10.17487/RFC5905.

\bibitem{elson2002} J. Elson, L. Girod, and D. Estrin, ``Fine-grained network time synchronization using reference broadcasts,'' in \emph{Proc. 5th Symp. Operating Systems Design and Implementation (OSDI)}, Boston, MA, USA, 2002; also \emph{ACM SIGOPS Operating Systems Review}, vol. 36, no. SI, pp. 147--163, 2002, doi: 10.1145/844128.844143.

\bibitem{olson2010} E. Olson, ``A passive solution to the sensor synchronization problem,'' in \emph{Proc. IEEE/RSJ Int. Conf. Intelligent Robots and Systems (IROS)}, 2010, pp. 1059--1064, doi: 10.1109/IROS.2010.5650579.

\bibitem{merry2019} K. Merry and P. Bettinger, ``Smartphone GPS accuracy study in an urban environment,'' \emph{PLOS ONE}, vol. 14, no. 7, Art. no. e0219890, 2019, doi: 10.1371/journal.pone.0219890.

\bibitem{shi2016} W. Shi, J. Cao, Q. Zhang, Y. Li, and L. Xu, ``Edge computing: Vision and challenges,'' \emph{IEEE Internet of Things Journal}, vol. 3, no. 5, pp. 637--646, 2016, doi: 10.1109/JIOT.2016.2579198.

\bibitem{barnell2018} M. Barnell, C. Raymond, C. Capraro, D. Isereau, C. Cicotta, and N. Stokes, ``High-performance computing (HPC) and machine learning demonstrated in flight using Agile Condor,'' in \emph{Proc. IEEE High Performance Extreme Computing Conf.}, 2018, pp. 1--4, doi: 10.1109/HPEC.2018.8547797.

\bibitem{perezgonzalez2025} A. P\'erez-Gonz\'alez, A. F. Villa-Salazar, I. N. Gomez-Miranda, J. D. Vel\'asquez-G\'omez, A. F. Romero-Maya, and A. Jaramillo-Duque, ``An IoV-based real-time telemetry and monitoring system for electric racing vehicles: Design, implementation, and field validation,'' \emph{Vehicles}, vol. 7, no. 4, Art. no. 128, 2025, doi: 10.3390/vehicles7040128.

\bibitem{liu2003} Y.-C. Liu, ``Effects of using head-up display in automobile context on attention demand and driving performance,'' \emph{Displays}, vol. 24, nos. 4--5, pp. 157--165, 2003, doi: 10.1016/j.displa.2004.01.001.

\bibitem{smith2016} M. Smith, J. L. Gabbard, and C. Conley, ``Head-up vs. head-down displays: Examining traditional methods of display assessment while driving,'' in \emph{Proc. 8th Int. Conf. Automotive User Interfaces and Interactive Vehicular Applications}, 2016, pp. 185--192, doi: 10.1145/3003715.3005419.

\bibitem{mareska2022} A. Mare\v ska, T. Kordov\'a, and M. Havl\'ik M\'ika, ``Production of Head-Up Display windshield and its relation with the image quality,'' \emph{Transactions on Transport Sciences}, vol. 13, no. 2, pp. 63--70, 2022, doi: 10.5507/tots.2022.011.

\bibitem{qin2017} Z. Qin, F.-C. Lin, Y.-P. Huang, and H.-P. D. Shieh, ``Maximal acceptable ghost images for designing a legible windshield-type vehicle Head-Up Display,'' \emph{IEEE Photonics Journal}, vol. 9, no. 6, pp. 1--12, 2017, doi: 10.1109/JPHOT.2017.2758820.

\bibitem{hauslschmid2016} R. H\"auslschmid, B. Pfleging, F. Alt, and A. Butz, ``A design space to support the development of windshield applications for the car,'' in \emph{Proc. 2016 CHI Conf. Human Factors in Computing Systems}, 2016, pp. 5076--5091, doi: 10.1145/2858036.2858336.

\bibitem{hauslschmid2015} R. H\"auslschmid, S. Osterwald, M. Lang, and A. Butz, ``Augmenting the driver's view with peripheral information on a windshield display,'' in \emph{Proc. 20th Int. Conf. Intelligent User Interfaces}, 2015, doi: 10.1145/2678025.2701393.

\bibitem{park2020} K. Park and Y. Im, ``Ergonomic guidelines of Head-Up Display user interface during semi-automated driving,'' \emph{Electronics}, vol. 9, no. 4, Art. no. 611, 2020, doi: 10.3390/electronics9040611.

\bibitem{reimer2014typeface} B. Reimer, B. Mehler, J. Dobres, J. F. Coughlin, S. Matteson, D. Gould, N. Chahine, and V. Levantovsky, ``Assessing the impact of typeface design in a text-rich automotive user interface,'' \emph{Ergonomics}, vol. 57, no. 11, pp. 1643--1658, 2014, doi: 10.1080/00140139.2014.940000.

\bibitem{you2021charsize} F. You, Y.-F. Yang, M.-T. Fu, J. Zhang, and J.-M. Wang, ``Design guidelines for the size and length of Chinese characters displayed in the intelligent vehicle's central console interface,'' \emph{Information}, vol. 12, no. 5, p. 213, 2021, doi: 10.3390/info12050213.

\bibitem{kim2014touchkey} H. Kim, S. Kwon, J. Heo, H. Lee, and M. K. Chung, ``The effect of touch-key size on the usability of in-vehicle information systems and driving safety during simulated driving,'' \emph{Applied Ergonomics}, vol. 45, no. 3, pp. 379--388, 2014, doi: 10.1016/j.apergo.2013.05.006.

\bibitem{euroncap2026engagement} Euro NCAP, \emph{Safe Driving: Driver Engagement Protocol}, version 1.2. Leuven, Belgium: Euro NCAP, 2026.

\bibitem{strayer2017infotainment} D. L. Strayer, J. M. Cooper, R. M. Goethe, M. M. McCarty, D. J. Getty, and F. Biondi, \emph{Visual and Cognitive Demands of Using In-Vehicle Infotainment Systems}. Washington, DC, USA: AAA Foundation for Traffic Safety, 2017.

\bibitem{mqtt311} A. Banks and R. Gupta, Eds., \emph{MQTT Version 3.1.1}, OASIS Standard, OASIS, 29 Oct. 2014. [Online]. Available: http://docs.oasis-open.org/mqtt/mqtt/v3.1.1/os/mqtt-v3.1.1-os.pdf

\bibitem{fall2003} K. Fall, ``A delay-tolerant network architecture for challenged internets,'' in \emph{Proc. ACM SIGCOMM Conf. Applications, Technologies, Architectures, and Protocols for Computer Communications}, Karlsruhe, Germany, 2003, pp. 27--34, doi: 10.1145/863955.863960.

\bibitem{chapungo2025} N. J. Chapungo and O. Postolache, ``Experimental evaluation of LoRaWAN connectivity reliability in remote rural areas of Mozambique,'' \emph{Sensors}, vol. 25, no. 19, Art. no. 6027, 2025, doi: 10.3390/s25196027.

\bibitem{ugwuanyi2021} S. Ugwuanyi, G. Paul, and J. Irvine, ``Survey of IoT for developing countries: Performance analysis of LoRaWAN and cellular NB-IoT networks,'' \emph{Electronics}, vol. 10, no. 18, Art. no. 2224, 2021, doi: 10.3390/electronics10182224.

\bibitem{liu2016ssd} W. Liu, D. Anguelov, D. Erhan, C. Szegedy, S. Reed, C.-Y. Fu, and A. C. Berg, ``SSD: Single shot multibox detector,'' in \emph{Computer Vision -- ECCV 2016}, Lecture Notes in Computer Science, vol. 9905. Cham, Switzerland: Springer, 2016, pp. 21--37, doi: 10.1007/978-3-319-46448-0\_2.

\bibitem{howard2017mobilenets} A. G. Howard, M. Zhu, B. Chen, D. Kalenichenko, W. Wang, T. Weyand, M. Andreetto, and H. Adam, ``MobileNets: Efficient convolutional neural networks for mobile vision applications,'' arXiv:1704.04861, 2017.

\bibitem{jacob2018quantization} B. Jacob, S. Kligys, B. Chen, M. Zhu, M. Tang, A. Howard, H. Adam, and D. Kalenichenko, ``Quantization and training of neural networks for efficient integer-arithmetic-only inference,'' in \emph{Proc. IEEE/CVF Conf. Computer Vision and Pattern Recognition}, 2018, pp. 2704--2713, doi: 10.1109/CVPR.2018.00286.

\bibitem{lin2014coco} T.-Y. Lin, M. Maire, S. Belongie, J. Hays, P. Perona, D. Ramanan, P. Doll\'ar, and C. L. Zitnick, ``Microsoft COCO: Common objects in context,'' in \emph{Computer Vision -- ECCV 2014}, Lecture Notes in Computer Science, vol. 8693. Cham, Switzerland: Springer, 2014, pp. 740--755, doi: 10.1007/978-3-319-10602-1\_48.

\bibitem{alqahtani2024edge} D. K. Alqahtani, M. A. Cheema, M. A. Rodriguez, and A. N. Toosi, ``A comprehensive evaluation of deep learning object detection models on heterogeneous edge devices,'' arXiv:2409.16808v3, Jul. 2026.

\bibitem{zhu2019distance} J. Zhu and Y. Fang, ``Learning object-specific distance from a monocular image,'' in \emph{Proc. IEEE/CVF Int. Conf. Computer Vision}, 2019, pp. 3838--3847, doi: 10.1109/ICCV.2019.00394.

\bibitem{lee1976tau} D. N. Lee, ``A theory of visual control of braking based on information about time-to-collision,'' \emph{Perception}, vol. 5, no. 4, pp. 437--459, 1976, doi: 10.1068/p050437.

\bibitem{dagan2004fcw} E. Dagan, O. Mano, G. P. Stein, and A. Shashua, ``Forward collision warning with a single camera,'' in \emph{Proc. IEEE Intelligent Vehicles Symposium}, 2004, pp. 37--42, doi: 10.1109/IVS.2004.1336352.

\bibitem{choe2023hazardnet} T. E. Choe, J. Wu, X. Lin, K. Kwon, and M. Park, ``HazardNet: Road Debris Detection by Augmentation of Synthetic Models,'' arXiv:2303.07547, 2023.

\bibitem{chan2021segmentmeifyoucan} R. Chan, K. Lis, S. Uhlemeyer, H. Blum, S. Honari, R. Siegwart, P. Fua, M. Salzmann, and M. Rottmann, ``SegmentMeIfYouCan: A Benchmark for Anomaly Segmentation,'' in \emph{Proc. Neural Information Processing Systems (NeurIPS) Track on Datasets and Benchmarks}, 2021.

\bibitem{tanveer2026yolov8n} A. Bint-e Tanveer, M. A. Kamal, M. M. Alam, and M. Mohd Su'ud, ``Real-time detection of rare roadside obstacles using YOLOv8-n in autonomous vehicles,'' \emph{PLOS ONE}, 2026, doi: 10.1371/journal.pone.0350732.

\bibitem{manglik2019} A. Manglik, X. Weng, E. Ohn-Bar, and K. M. Kitani, ``Forecasting time-to-collision from monocular video: Feasibility, dataset, and challenges,'' in \emph{Proc. IEEE/RSJ Int. Conf. Intelligent Robots and Systems (IROS)}, 2019, pp. 8081--8088, doi: 10.1109/IROS40897.2019.8967730.

\bibitem{lenkeit2012fcw} J. Lenkeit, ``NCAP forward collision warning confirmation test: 2013 Ford Explorer,'' Dynamic Research, Inc., report OCAS-DRI-FCW-13-009, prepared for the National Highway Traffic Safety Administration, Washington, DC, USA, Oct. 2012. [Online]. Available: https://nrd-static.nhtsa.dot.gov/reports/vehdb/v00000/v08900/v08966R001.pdf

\bibitem{desimone2018} M. De Simone, Z. B. Rivera Chavez, and D. Guida, ``Obstacle Avoidance System for Unmanned Ground Vehicles by Using Ultrasonic Sensors,'' \emph{Machines}, vol. 6, no. 2, Art. no. 18, 2018, doi: 10.3390/machines6020018.

\bibitem{noomwongs2025} N. Noomwongs, K. Kiataramgul, S. Chantranuwathana, and G. Phanomchoeng, ``GNSS-RTK-Based Navigation with Real-Time Obstacle Avoidance for Low-Speed Micro Electric Vehicles,'' \emph{Machines}, vol. 13, no. 6, Art. no. 471, 2025, doi: 10.3390/machines13060471.

\bibitem{caroleo2026} G. Caroleo, A. Albini, and P. Maiolino, ``On the Characterisation of the Time-of-Flight VL53L5CX Sensor by STMicroelectronics for Indoor Robotics Applications,'' \emph{Sensors}, vol. 26, no. 5, Art. no. 1639, 2026, doi: 10.3390/s26051639.

\bibitem{yeong2021} D. J. Yeong, G. Velasco-Hernandez, J. Barry, and J. Walsh, ``Sensor and sensor fusion technology in autonomous vehicles: A review,'' \emph{Sensors}, vol. 21, no. 6, Art. no. 2140, 2021, doi: 10.3390/s21062140.

\bibitem{eising2022} C. Eising, J. Horgan, and S. Yogamani, ``Near-field perception for low-speed vehicle automation using surround-view fisheye cameras,'' \emph{IEEE Transactions on Intelligent Transportation Systems}, vol. 23, no. 9, pp. 13976--13993, 2022, doi: 10.1109/TITS.2021.3127646.

\bibitem{rimell2007} A. N. Rimell and N. J. Mansfield, ``Design of digital filters for frequency weightings required for risk assessments of workers exposed to vibration,'' \emph{Industrial Health}, vol. 45, no. 4, pp. 512--519, 2007, doi: 10.2486/indhealth.45.512.

\bibitem{virtins2024} Virtins Technology, ``Frequency weightings for sound and vibration perceived by humans using Multi-Instrument,'' application note, 2024. [Online]. Available: \url{https://virtins.com/doc/Frequency-Weightings-for-Sound-and-Vibration-Perceived-by-Humans-using-Multi-Instrument.pdf} (accessed Oct. 24, 2026).

\bibitem{makerfabsa7670x} Makerfabs, ``ESP32S3 4G LTE CAT1 A7670X,'' product page and hardware repository. [Online]. Available: \url{https://www.makerfabs.com/esp32s3-4g-lte.html} and \url{https://github.com/Makerfabs/ESP32-S3-4G-LTE-CAT1-A7670X} (accessed 24 Sep. 2026).

\bibitem{invensense2012rm} InvenSense Inc., ``MPU-6000 and MPU-6050 register map and descriptions,'' Document RM-MPU-6000A-00, Rev.~4.0, Sunnyvale, CA, USA, Mar. 2012.

\bibitem{agostinacchio2014} M. Agostinacchio, D. Ciampa, and S. Olita, ``The vibrations induced by surface irregularities in road pavements: A Matlab approach,'' \emph{European Transport Research Review}, vol. 6, pp. 267--275, 2014, doi: 10.1007/s12544-013-0127-8.

\bibitem{simcoma76xx} SIMCom Wireless Solutions, \emph{A76XX Series AT Command Manual}, V1.09, Apr. 2023. [Online]. Available: \url{https://files.waveshare.com/wiki/A7670E-Cat-1-GNSS-HAT/A76XX_Series_AT_Command_Manual_V1.09.pdf} (accessed 24 Sep. 2026).

\bibitem{sensirionsen5x} Sensirion AG, \emph{Datasheet SEN5x: Environmental Sensor Node for HVAC and Air Quality Applications}, version 2, Stäfa, Switzerland, Mar. 2022.

\bibitem{picamera2manual} Raspberry Pi Ltd, \emph{The Picamera2 Library}, Sep. 2026. [Online]. Available: \url{https://datasheets.raspberrypi.com/camera/picamera2-manual.pdf} (accessed 28 Sep. 2026).

\bibitem{litert} Google, ``LiteRT,'' Google AI Edge documentation. [Online]. Available: \url{https://ai.google.dev/edge/litert}

\bibitem{ultralytics2024yolo11} G. Jocher and J. Qiu, \emph{Ultralytics YOLO11}, version 11.0.0, software, 2024. [Online]. Available: \url{https://github.com/ultralytics/ultralytics}

\bibitem{gerapothole} gerapothole, ``pothole-detection-yolov8,'' dataset, version 1, Roboflow Universe. [Online]. Available: \url{https://universe.roboflow.com/gerapothole/pothole-detection-yolov8} (accessed 2 Oct. 2026).

\bibitem{nienaber2015dataset} S. Nienaber, M. J. Booysen, and R. S. Kroon, ``Detecting potholes using simple image processing techniques and real-world footage,'' in \emph{Proc. 34th Southern African Transport Conf.}, Pretoria, South Africa, 2015. Image dataset redistributed as ``Road pothole images for pothole detection,'' Kaggle (sovitrath). [Online]. Available: \url{https://www.kaggle.com/datasets/sovitrath/road-pothole-images-for-pothole-detection}

\bibitem{ncnn} Tencent, \emph{ncnn: A High-Performance Neural Network Inference Framework Optimized for the Mobile Platform}, software. [Online]. Available: \url{https://github.com/Tencent/ncnn}

\bibitem{sen1968} P. K. Sen, ``Estimates of the regression coefficient based on Kendall's tau,'' \emph{Journal of the American Statistical Association}, vol. 63, no. 324, pp. 1379--1389, 1968, doi: 10.1080/01621459.1968.10480934.

\bibitem{saej3016} SAE International, \emph{Taxonomy and Definitions for Terms Related to Driving Automation Systems for On-Road Motor Vehicles}, SAE J3016, Apr. 2021.

\bibitem{iso26262} International Organization for Standardization, \emph{Road Vehicles: Functional Safety}, ISO 26262:2018.

\bibitem{iso21448} International Organization for Standardization, \emph{Road Vehicles: Safety of the Intended Functionality}, ISO 21448:2022.

\bibitem{raspberrypigpio} Raspberry Pi Ltd, ``Raspberry Pi hardware: GPIO and the 40-pin header,'' Raspberry Pi documentation. [Online]. Available: \url{https://www.raspberrypi.com/documentation/computers/raspberry-pi.html}

\bibitem{luxen2011osrm} D. Luxen and C. Vetter, ``Real-time routing with OpenStreetMap data,'' in \emph{Proc. 19th ACM SIGSPATIAL Int. Conf. Advances in Geographic Information Systems}, 2011, pp. 513--516, doi: 10.1145/2093973.2094062.

\bibitem{osm} OpenStreetMap contributors, \emph{OpenStreetMap}, map data available under the Open Database License. [Online]. Available: \url{https://www.openstreetmap.org}

\bibitem{rimell2010corr} A. N. Rimell and N. J. Mansfield, ``Corrigendum: Design of digital filters for frequency weightings required for risk assessments of workers exposed to vibration,'' \emph{Industrial Health}, vol. 48, no. 5, p. 743, 2010.

\end{thebibliography}
